\section{Related Work}
\label{sec:related_work}

\subsection{Sampling-based \acs*{MPC}}

A prominent approach to sampling-based \ac{MPC} is the \ac{CEM}, which iteratively updates a proposal \ac{PDF} by selecting low-cost elite samples and maximizing their likelihood~\cite{rubinsteinCrossEntropy2004}.
For Gaussian proposals, the updated moments for the next iteration are computed as the empirical mean and covariance of the selected control sequences~\cite{rubinsteinCrossEntropy2004}.
To improve numerical robustness and reduce premature concentration of the proposal, many implementations introduce a momentum (or smoothing) term across iterations~\cite{rubinsteinCrossEntropy2004,pinneriSampleefficientCrossentropyMethod2021}.
For the mean update, this is typically written as
\begin{align}
    \label{eq:cem_momentum_mean}
    \evVarCem_{j+1} = \alpha \, \evVarCem_{j} + (1-\alpha) \, \evVarCem_{j}^{\prime} \enspace,
\end{align}
where $\evVarCem_{j}^{\prime}$ denotes the (unsmoothed) mean of the elite set.

Related \ac{MPPI}-like methods also perform (iterative) proposal refinement, but they rely on importance weighting of sampled control sequences instead of hard elite selection~\cite{arXiv26_Walker_IFAC}.
The weights are given by
\begin{align}
    \label{eq:imp_weights}
    w^{(i)} \propto \exp\left(-\frac{1}{\lambda} J\left(\vVarCem^{(i)}\right) \right) \enspace,
\end{align}
and are used to compute the updated Gaussian moments as the control sequences' weighted sample mean and covariance.
The inverse temperature $\lambda$ controls the exploration-exploitation trade-off and remains constant~\cite{williamsInformationtheoretic2018} or is adapted heuristically during optimization~\cite{pezzatoSamplingbasedModelPredictive2025}.
% Following~\cite{pezzatoSamplingbasedModelPredictive2025}, one representative adaptation is
% \begin{align}
%     \label{eq:temp_adaption}
%     \lambda_{j+1} = \begin{cases} 
%         0.9 \, \lambda_j & \text{if~} \eta > \eta_{\mathrm{max}}
%         \\
%         1.2 \, \lambda_j & \text{if~} \eta < \eta_{\mathrm{min}}
%         \\
%         \lambda_j & \text{otherwise}        
%         \end{cases} \enspace,
% \end{align}
% where $\eta$ is the normalization constant from \cref{eq:imp_weights},
% %, which serves as a proxy for the effective sample size.
% and $\lambda_j$ and $\lambda_{j+1}$ are the inverse temperatures in iterations $j$ and $j+1$, respectively.
% Thresholds $\eta_{\mathrm{min}}$ and $\eta_{\mathrm{max}}$ are selected empirically to ensure smooth behavior; \eg,~\cite{pezzatoSamplingbasedModelPredictive2025} reports good performance for the range $[5,10]$.
Furthermore, iterative \ac{MPPI} variants additionally employ momentum-style smoothing~\cref{eq:cem_momentum_mean} of the proposal update, analogous to \ac{CEM}, to mitigate premature convergence~\cite{bhardwajSTORMIntegratedFramework2022}.

% Improvements in sampling
Orthogonal to the above methods, improvements in sampling efficiency and smoothness of the optimized control inputs have been achieved by using low-discrepancy sequences, such as the Halton sequence~\cite{bhardwajSTORMIntegratedFramework2022}, time-correlated noise~\cite{pinneriSampleefficientCrossentropyMethod2021}, or, as in our previous work, by using optimal deterministic samples based on the \ac{LCD}~\cite{walkerSampleefficientSmoothCrossentropy2026,arXiv26_Walker_IFAC}.


\subsection{Trust Regions in \ac*{RL}}
In \ac{RL}, trust region methods provide principled policy updates that balance exploration and exploitation while ensuring stable learning.
The core idea is to limit the update step by constraining the \ac{KL} divergence between the new policy and the old policy, preventing drastic changes that can lead to performance degradation.
Adding an entropy constraint further encourages exploration and prevents premature convergence to suboptimal policies.
These concepts have been established for policy search~\cite{petersRelativeEntropyPolicy2010,abdolmaleki2015model}, step-based \ac{RL}~\cite{schulmanTrustRegionPolicy2015}, episodic \ac{RL}~\cite{ottoDeepBlackboxReinforcement2023}, and gradient-based stochastic optimal control with quadratic reward functions~\cite{blessingTrustRegionConstrained2025}.